\section{Main Result}

\subsection{Periodic target and current status}

This draft does not currently prove the periodic theorem.  The checked open
ingredient is the signed critical-height sign-persistence / peak-height theorem behind the Gold \(L^1\) source wall.  The statement below is
the continuation consequence that would follow once that coupled storage object
is proved from the original parent packet before child clipping.

\begin{proposition}[Conditional periodic continuation consequence]
Let \(\nu>0\). Let
\[
u_0\in C^\infty(\mathbb T^3;\mathbb R^3),\qquad
\nabla\cdot u_0=0,\qquad \int_{\mathbb T^3}u_0(x)\,dx=0.
\]
Then the continuation criterion for the incompressible Navier--Stokes system
\[
\partial_tu+(u\cdot\nabla)u+\nabla p=\nu\Delta u,
\qquad \nabla\cdot u=0,
\qquad u(0)=u_0,
\]
would force the local periodic solution to extend smoothly through every finite
time interval:
\[u\in C^\infty_{\mathrm{loc}}([0,\infty)\times\mathbb T^3),
\qquad p\in C^\infty_{\mathrm{loc}}([0,\infty)\times\mathbb T^3),
\]
with the usual periodic pressure normalization.
\end{proposition}

\subsection{Four standard lemmas}

\begin{lemma}[Moving-cylinder regularity and affine excess]
Let \(Q_r^\Phi\) be a same-fluid moving parabolic cylinder with bounded distortion. If its normalized scale-critical CKN packet is below a fixed threshold \(\varepsilon_m\), then a smaller cylinder carries finite-depth derivative and pressure bounds. On that smaller cylinder the affine-subtracted velocity residual satisfies
\[
D^+X_R+cN_R\le B_RX_R+C X_R^{1/2}N_R+F_R.
\]
On scheduler intervals the nonlinear term is absorbed, and
\[
X_R\in L^\infty,\qquad N_R\in L^1,
\qquad \mathcal K_{\le m}^{avg,R}\in L^1.
\]
\end{lemma}

\begin{lemma}[Terminal good-cover dichotomy and globalization]
For a bounded-distortion same-fluid terminal tail, exactly one of the following alternatives holds. Either the tail has a terminal positive-scale family of scale-critical good cylinders, or every terminal subtail loses all uniform positive good scale. In the first case, a finite cover with bounded overlap and a common scheduler globalizes the preceding local estimates into a terminal averaged transported-tower bound.
\end{lemma}

\begin{lemma}[Conditional endpoint contradiction]
A first finite terminal endpoint must enter one of four primitive faces: participation failure, packing failure, averaged tower-amplitude escape, or terminal good-scale loss. The terminal package eliminates these four faces respectively by participation, packing, averaged tower control, and terminal field-scale persistence. Therefore no first finite endpoint occurs.
\end{lemma}

\begin{lemma}[Terminal readout and continuation]
The averaged endpoint package gives a finite positive-scale readout cover. On that cover, the averaged field and tower bounds recover the pointwise endpoint package. The pointwise package gives
\[
\sup_{t<T_*}\|u(t)\|_{H^s(\mathbb T^3)}<\infty
\]
for some \(s>5/2\). The standard periodic \(H^s\) local theory then extends the solution past any alleged finite endpoint.
\end{lemma}

\subsection{Conditional proof of the continuation consequence}

Let \((u,p)\) be the maximal classical solution on \([0,T_*)\). Suppose \(T_*<\infty\). By the terminal good-cover dichotomy, the terminal same-fluid tail either admits a positive-scale good-cylinder cover or loses all uniform positive good scale. On the positive-scale branch, the moving-cylinder regularity and affine-excess lemma globalizes to the terminal transported-tower bound. On the scale-loss branch, the loss is one of the endpoint faces.

The endpoint contradiction lemma excludes every possible first finite endpoint face. Hence the terminal endpoint package holds. The terminal readout and continuation lemma then gives a uniform \(H^s(\mathbb T^3)\) bound, with \(s>5/2\), and the periodic local theory extends the solution beyond \(T_*\). This contradicts maximality. Therefore \(T_*=\infty\).

Smoothness on every finite time interval follows from the standard smooth-data periodic theory. The solution is global and smooth.

\paragraph{Scope.}
The theorem is periodic on \(\mathbb T^3\). A whole-space theorem on \(\mathbb R^3\) is a separate export problem.
